% 图4-2：一次模型标注中的四类输入，按112 mm正文宽度设计。
\begin{tikzpicture}[
  x=1mm,
  y=1mm,
  font=\figurefont\footnotesize,
  text=BookInk,
  input block/.style={book node,draw=FigureInput,fill=FigureInput!6,line width=0.8pt,text width=37mm,minimum height=12mm,inner sep=1.4mm},
  model block/.style={book node,draw=FigureModel,fill=FigureModel!8,line width=0.9pt,text width=25mm,minimum height=18mm,inner sep=1.5mm},
  output block/.style={book node,draw=FigureOutput,fill=FigureOutput!6,line width=0.8pt,text width=17mm,minimum height=18mm,inner sep=1.2mm},
  convention/.style={-{Stealth[length=2mm]},draw=BookAmber,line width=0.8pt,dashed},
  evidence/.style={-{Stealth[length=2mm]},draw=BookBlue,line width=0.9pt},
  edge text/.style={fill=white,inner sep=0.7pt,font=\figurefont\scriptsize,text=BookMuted}
]
  \node[input block,draw=BookAmber] (requirement) at (21,0) {任务要求与输出约定\\物种表、弃权、返回字段};
  \node[input block,draw=BookAmber] (demonstration) at (21,-18) {示范案例与参照答案\\S011：山茶；S017：弃权};
  \node[input block] (sample) at (21,-36) {待处理样本与必要上下文\\S042：花、叶、地点、时间};
  \node[input block] (support) at (21,-54) {辅助证据与关联信息\\本地名录R3、带标签近邻};

  \node[model block] (model) at (70,-27) {标注模型\\版本$M_5$};
  \node[output block] (output) at (101,-27) {候选记录\\答案、分数\\证据与版本};

  \draw[convention] (requirement.east) -- node[edge text,above] {约定} ([yshift=6mm]model.west);
  \draw[convention] (demonstration.east) -- node[edge text,above] {示范} ([yshift=2mm]model.west);
  \draw[evidence] (sample.east) -- node[edge text,below] {当前观测} ([yshift=-2mm]model.west);
  \draw[evidence] (support.east) -- node[edge text,below] {外部证据} ([yshift=-6mm]model.west);
  \draw[-{Stealth[length=2mm]},draw=FigureModel,line width=1.0pt] (model.east) -- (output.west);
\end{tikzpicture}
